\documentclass[hidelinks]{aquila-doc}

% For removing page numbers
% \pagenumbering{gobble}

% Aquila does not come with a built-in reference manager
% You can choose your own package and here we uses biblatex
\usepackage{biblatex}
\addbibresource{sources.bib}
% If you are partial to natbib you can use it too
% \usepackage[numbers]{natbib}
% \bibliographystyle{plainnat}

\title{\aquila is a modern template for \LaTeX documents}
\author{}
\date{}

\begin{document}

\maketitle

While you can still write text across the entire document, \aquila uses the \command{goldencolumn} command to write content in asymmetric columns by default.
The right column is widest and used for main content.
The left column can be used for any additional content that is not integral to the main story, but can clarify the exposition.

\goldencolumn{
    \section{Overview of the basic principles of an \aquila document}

    You have now entered the first section of the main story of the document.
    All content should be placed inside a \command{goldencolumn}, which includes every \command{section} and \command{subsection}.

    \aquila promotes the use of whitespace for a clearer and cleaner exposition.
    Whitespace becomes available by leaving open a smaller left column that can be used at the writer's leisure.
    The column can give additional insights, clarifications, comments etc.
    These additions are easily added through the \command{comment}.

    \comment{
    Content in the \command{comment} environment will always align with the content that follows right after the environment.
    }

    In other words, a \command{comment} appears right before the text you want to comment.\footnote{Footnotes should be used sparingly as most of their functionality should be covered by comments.}

    Mathematics in \aquila uses sans-serif fonts just like any other text.
    For example, the mathematical facts that
    \begin{align}
        \int_{-\infty}^{\infty}\
        e^{-x^2}
        \diff x
        =
        \sqrt\pi
        \quad
        \text{and}
        \quad
        e^{i \pi} + 1 = 0,
    \end{align}
    which illustrate the relationship between $\pi$ and $e$.

    While the wider default column of \aquila provides more space for mathematical expressions than a traditional two-column layout, it is still sometimes insufficient.
    For equations that are unavoidably longer or deserve central attention, you can use the \command{spanalign} command instead, e.g. the beautiful Einstein field equations~\cite{einstein1916}
    \spanalign{
        G_{\mu \nu} + \Lambda g_{\mu \nu}
        =
        \frac{8\pi G}{c^4} T_{\mu \nu}
    }
    For regular pictures and figures that should span across columns, you can use the starred version of the \texttt{figure} environment.

    \subsection{Tikz can be a lot of fun}

    \comment{\shift{Overview}}

    Aquila promotes using \texttt{tikz} as a language for clean graphics.
    Do to so, it provides a number of default node types and arrow styles, as well as useful commands for pointing at text.
    These macros, settings and commands should make using \texttt{tikz} at least a little less cumbersome as well as provide consistent graphical representations of your figures.
    One simple example is a \accent{Bayesian network}.

    \comment{
    \begin{tikzpicture}
        \node[circnode] (b) {\texttt{b}};
        \node[circnode, right=3em of b.east] (e) {\texttt{e}};
        \node[between=b and e] (middle) {};
        \node[circnode, below=2em of middle] (a) {\texttt{a}};

        \draw[thinpath] (b.south) |- (a.west);
        \draw[thinpath] (e.south) |- (a.east);
    \end{tikzpicture}
    }

    \begin{example}[Bayesian network]
        Let \texttt{b}, \texttt{e} and \texttt{a} be three random variables that represent the occurrence of a burglary, earthquake or alarm, respectively.
        Their dependencies are represented in the probabilistic graphical model (PGM) on the left.
    \end{example}

    The \texttt{tikz} code of this example shows a number of frequently used macros.
    It uses \texttt{circnode}'s as circular nodes traditional in a PGM as well as \texttt{thinpath}'s to draw the dependencies.
    Placing such nodes and arrows can quickly become convoluted or unorganised, so we recommend using the \texttt{between} macro as well as only relative distances to define intermediate nodes.
    This practice should help maintain consistent alignment as well as overview.
    Specificying cardinal directions, \eg \texttt{east} or \texttt{south}, is also recommended.

    Apart from graphical models or flowcharts, \texttt{tikz} is also used in aquila to \accent{connect parts of mathematics to their textual representation}.
    These connections are situational and don't always work, but can help with clarity of exposition.

    \comment{
    Where \colouraccentnode{orange_red}{recursive samples}{recursion_txt} are \colouraccentnode{celadon_blue}{transitioned}{transition_txt} and then \colouraccentnode{jungle_green}{reweighed by observations}{observation_txt}.
    }

    For instance, consider the recursive particle filtering equation
    \begin{align}
        \int
        \colourmathnode{jungle_green}{
            p(\vector{Z}_{t + 1} \mid \vector{X}_{t + 1})
        }{observation}
        \cdot
        \colourmathnode{celadon_blue}{
            p(\vector{X}_{t + 1} \mid \vector{x}_{t})
        }{transition}
        \cdot
        \colourmathnode{orange_red}{
            p(\vector{x}_t \mid \vector{Z}_{0:t})
        }{recursion}
        \diff\vector{x}_t
    \end{align}
    The two most important macros here are \command{colouraccentnode} and \command{colourmathnode} that turn either a piece of accented text or a piece of mathematics into a \texttt{tikz} node with a given name.
    These nodes can then be connected in an \texttt{overlaypicture} if one desires, although utilising only colour can already be clarifying.
    While connecting text and mathematics does require some finetuning, the result can be quite pleasing.
    In case connections are not desired, then one can simply use \command{colouraccent} or \command{colourmath} instead.

    \begin{overlaypicture}
        \draw[thinline, jungle_green]
        (observation_txt) |-
        ([yshift=-0.75em]observation.south) --
        (observation.south);
        \draw[thinline, celadon_blue]
        (transition_txt.east) --
        ([xshift=0.75em]transition_txt.east) |-
        ([yshift=1em]transition.north) --
        (transition.north);
        \draw[thinline, orange_red]
        (recursion_txt.north) |-
        ([yshift=1em, xshift=-8em]recursion_txt.north) |-
        ([yshift=-1em]recursion.south) --
        (recursion.south);
    \end{overlaypicture}

    \subsection{Other nice things to know}

    Referencing a series of words can be quite useful.
    Say for instance that you have the following research question.

    \comment{\shift{Research question I\namelabel{Research question I}}}

    \begin{flushright}
        \qemph{How to write a clean document?}

    \end{flushright}
    If you like to refer to this question with a full link, then you can do so my just writing that \ref{Research question I} is quite important.

    Other things will be added here in the future.

    \section{Section headers describe the section content in a short sentence}

    Before we end this example, a short note on how to generally approach writing in \aquila.
    When writing, each paragraph should start with a topical sentence roughly describing the content of the paragraph.
    Each subsequent sentence then goes into more detail by either going deeper into the object of the previous sentence or adding more information to the object of the previous sentence.
    After these sentences, you can then get to the main point of the paragraph again if applicable.

    \paragraph{Very important paragraphs can be emphasised by the \command{paragraph} command to give them more visibility.}
    Sometimes a point is so important that you really want to make sure people remember or that it is noticed even when skimming through the document.
    For that, \command{paragraph} puts the topical sentence in bold.
}

% I am partial to not having references directly follow the text
\newpage

% For the default biblatex
\printbibliography
% For natbib
% \bibliography{sources}

\newpage
\appendix

\label{app:test}

\goldencolumn{
    \section{The appendix contains the juicy details of your work that didn't make the cut}

    This part is where the fun really begins by explaining all the gritty details of your work that were too obscure to be entertained in the main story.
}

% \bibliography{sources}

\end{document}
